% ============================================================================
% Generic manuscript scaffold — copy this directory to start a new project.
% Class options below target Phys. Rev. A by default; swap the commented
% alternatives for PRApplied / PRX Quantum / PRL / a generic preprint.
% ============================================================================
\documentclass[aps,pra,reprint,superscriptaddress,longbibliography]{revtex4-2}
% \documentclass[aps,prapplied,reprint,superscriptaddress,longbibliography]{revtex4-2}   % Phys. Rev. Applied
% \documentclass[aps,prxquantum,reprint,superscriptaddress,longbibliography]{revtex4-2}  % PRX Quantum
% \documentclass[aps,prl,twocolumn,superscriptaddress]{revtex4-2}                        % Phys. Rev. Lett.

\usepackage{graphicx}
\usepackage{amsmath,amssymb}
\usepackage{bm}
\usepackage{hyperref}

% Shared project conventions (keep consistent with
% skills/literature-and-notation/references/notation.md):
%   natural units hbar^2/2m_e = 1, hbar = 1
\newcommand{\ehe}{electrons on helium}

\begin{document}

\title{TITLE}

\author{First Author}
\author{Second Author}
\affiliation{Affiliation}

\date{\today}

\begin{abstract}
One-paragraph abstract: state the problem, the idea, and (once available) the headline result.
Keep it hedged appropriately if the manuscript is still a proposal rather than a set of results.
\end{abstract}

\maketitle

\section{Introduction}
Motivate the problem. Ground it in what has already been demonstrated or established (cite prior
work by BibTeX key), then state the gap this project addresses.

\section{Model}
Define the Hamiltonian / framework. State units and conventions explicitly (point to
\texttt{notation.md} rather than re-deriving the project's shared conventions from scratch).

\section{Results}
Placeholder — fill in once computed. Distinguish demonstrated/derived results from
projected/simulated ones.

\section{Discussion and outlook}
Placeholder.

\begin{acknowledgments}
Placeholder.
\end{acknowledgments}

\appendix
\section{Appendix title}
Placeholder for supporting derivations.

\bibliography{references}

\end{document}
